\chapter*{\chapterstyle{III --- Introduction à la réduction}}
\addcontentsline{toc}{section}{Introduction à la réduction}
Dans ce chapitre, nous étudirons un domaine vaste de l'algèbre linéaire appellé \textbf{réduction des endormorphismes}.\+
En effet, sous une forme quelconque un endomorphisme représenté par une matrice présente plusieurs problèmes:
\begin{align*}
   &\bullet \;\; \text{Il est couteux de calculer les puissances d'une matrice quelconque.} \\
   &\bullet \;\; \text{Une représentation quelconque donne peu d'informations sur l'endomorphisme.}
\end{align*}
L'objectif sera donc de réduire (comprendre simplifier) la représentation de l'endormorphisme, et de le représenter par une matrice plus simple. \<

Soit \(f \in \mathcal{L}(E)\) et une famille \((E_i)\) de sous-espaces \textbf{supplémentaires}\footnote[1]{Comme tout sous-espace admet un supplémentaire (théorème de la base incomplète), on peut définir une base adaptée à \textbf{un seul} sous-espace comme étant une base adaptée à la somme directe de ce sous-espace et de son supplémentaire, qui revient à compléter la base en une base de \(E\).}, alors on peut construire une base \(\mathcal{B}\) dite \textbf{base adaptée à la décomposition} en concaténant des bases respectives des \(E_i\).

\subsection*{\subsecstyle{Elements Propres {:}}}
Soit \(\lambda \in \R\), on dit que \(\lambda\) est une \textbf{valeur propre}\footnote[2]{On appelle \textbf{specte}, noté Sp\((f)\) l'ensemble des valeurs propres de \(f\).} de l'endomorphisme \(f\) si et seulement si il existe un vecteur non-nul \(u \in E\) tel que:
\customBox{width=2.5cm}{
   \(f(u) = \lambda u\)
}
On dira alors qu'un tel vecteur est \textbf{vecteur propre} de l'endomorphisme et on appelle \textbf{sous-espace propre} associé à la valeur propre \(\lambda\) l'ensemble des vecteurs propres associés, qu'on note \(E_\lambda\) dont on déduit une expression\footnote[3]{Directement d'après la définition d'un vecteur propre associé à \(\lambda\).}:
\customBox{width=5cm}{
   \(E_\lambda = \text{Ker}(f - \lambda \text{Id})\)
}
Enfin, on peut montrer une propriété trés importante pour la suite:
\customBox{width=10cm}{
   Toute somme de sous-espaces propres est \textbf{directe}.
}
\begin{center}
   \textit{L'image des sous-espaces propres par l'endomorphisme se réduit à une homotéthie de rapport la valeur propre.}
\end{center}

\subsection*{\subsecstyle{Sous-Espaces Stables {:}}}
On dit que \(F\) est \textbf{stable par l'endormorphisme} si et seulement si:
\customBox{width=2.5cm}{
   \(f(F) \subseteq F\)
}
On considère maintant une base de \(F\) qu'on compléte en une base de \(E\) via le théorème de la base incomplète, alors dans une telle base, l'endormorphisme est représenté par la matrice par blocs:
\[
   \left(\begin{array}{c|c}
      A & B\\
      \hline\\[-1.7\medskipamount]
      0 & C
   \end{array}\right)
\]
\pagebreak

Plus généralement, si on a une famille \((E_i)\) de sous-espaces \textbf{stables et supplémentaires}, ie \(E = \bigoplus_{i \in \N} E_i\), et \(\mathcal{B}\) est une base adaptée à cette décomposition, alors dans cette base, \(f\) est représentée par la matrice diagonale par blocs:
\[
   \left(\begin{array}{ccc}
      A_1 & {} & {}\\
      {} & \ddots & {}\\
      {} & {} & A_n\\
   \end{array}\right)
\]
On remarque donc que la stabilité des sous-espaces nous permet de reprénsenter notre transformation de manière plus simple, en particulier, on peut alors montrer une propriété fondamentale:
\customBox{width=10cm}{
   \textbf{Tout les sous-espaces propres sont stables.}
}

\subsection*{\subsecstyle{Polynôme caractéristique {:}}}
On peut montrer\footnote[1]{\(E_\lambda\) est un noyau, il suffit de caractériser le fait qu'il soit non vide en termes de la bijectivité d'un certain endomorphisme.} que \(\lambda\) est valeur propre si et seulement si:
\begin{align*}
   E_\lambda \neq \bigl\{0_E\bigl\} \Longleftrightarrow \text{det}(f - \lambda \text{Id}) = 0
\end{align*}
On définit alors le \textbf{polynôme caractéristique} d'un endormorphisme par:
\customBox{width=5cm}{
   \(P_f=\text{det}(f - X \text{Id})\)
}
En particulier, on peut donc montrer que:
\customBox{width=15cm}{
   \textbf{Les valeurs propres sont exactement les racines du polynôme caractéristique.}
}
Ceci nous donne donc une méthode systématique pour trouver les valeurs propres d'un endomorphisme. En particulier, on peut alors montrer les identités suivantes, utiles dans la recherche de valeurs propres:
\begin{align*}
   &\bullet \;\; \text{La somme des valeurs propres est égale à \textbf{la trace de la matrice}.} \\
   &\bullet \;\; \text{Le produit des valeurs propres est égale au \textbf{déterminant de la matrice}.}
\end{align*}

\subsection*{\subsecstyle{Diagonalisation {:}}}
Les endomorphismes qu'on peut représenter le plus simplement sont ceux qui ceux réduisent (dans une base bien choisie) à une homotéthie des vecteurs de la base. On dira alors que ces endomorphismes sont \textbf{diagonalisables}, formellement:
\customBox{width=16.5cm}{
   Un endormorphisme est diagonalisable si il existe une base de \(E\) constituée de vecteurs propres de \(f\).
}
De manière équivalente:
\customBox{width=16cm}{
   Un endormorphisme est diagonalisable si ses matrices sont semblables à une matrice diagonale.
}

\underline{Exemple:} Soit \(f\) un tel endormorphisme de \(\R^3\) et \((e_{\lambda_1}, e_{\lambda_2}, e_{\lambda_3})\) des tels vecteurs propres, alors dans cette base, \(f\) est représenté par la matrice:
\[
   D = \left(\begin{array}{ccc}
      \lambda_1 & {} & {}\\
      {} & \lambda_2 & {}\\
      {} & {} & \lambda_3
   \end{array}\right) 
\]
Ou encore si \(A\) est la matrice de \(f\) dans la base canonique, alors \(A = PDP^{-1}\) avec:
\[
   P = ([e_{\lambda_1}]_\mathscr{C}, [e_{\lambda_2}]_\mathscr{C}, [e_{\lambda_3}]_\mathscr{C}) \text{ et } D = \left(\begin{array}{ccc}
      \lambda_1 & {} & {}\\
      {} & \lambda_2 & {}\\
      {} & {} & \lambda_3
   \end{array}\right) 
\]
\begin{center}
   \textit{Diagonaliser un endomorphisme revient à \textbf{décomposer l'espace en somme directe de droites stables}.
   }
\end{center}
\pagebreak

\subsection*{\subsecstyle{Critères de diagonalisabilité {:}}}
On sait qu'un endomorphisme \(f\) est diagonalisable si et seulement si il admet une base de vecteurs propres, alors on peut montrer\footnote[1]{En effet, si tel est le cas, alors il suffit de prendre une base pour chaque \(E_\lambda\) (qui est bien constituée de vecteurs propres par définition), et de les concaténer pour obtenir une base de \(E\) constituée de vecteurs propres.} que \(f\) est diagonalisable si et seulement si:
\customBox{width=4cm}{
   \(
      E = \bigoplus_{\lambda \in \text{Sp}(f)} E_\lambda  
   \)
}
En particulier on remarque alors que montrer la supplémentarité revient à montrer que \textbf{la somme des dimensions des sous-espaces propres est égale à la dimension totale} car toute somme de sous-espaces propres est directe.\<

On peut alors montrer que si \(\lambda\) est une valeur propre de multiplicité \(\alpha\) pour le polynome caractéristique, alors:
\[
   1 \leq \text{dim}(E_\lambda) \leq \alpha   
\]
On a alors le théorème fondamental suivant:
\begin{center}
   \textbf{Un endomorphisme est diagonalisable si et seulement si son polynôme est scindé sur \(\K\) et que la dimension de chaque sous-espace propre est égale à la multiplicité de la valeur propre associée}.
\end{center}
On peut donc étudier si un endormorphisme est diagonalisable en calculant les valeurs propres et les dimensions des sous-espaces propres associés, ce qui revient à un calcul de polynôme caractéristique suivi de calculs de noyau.

\underline{Exemple:} Diagonalisons la matrice \(A = \left(\begin{array}{ccc}
   1 & 1 & 1\\
   2 & 2 & 2\\
   3 & 3 & 3
\end{array}\right) \)\< 

On peut calculer \(P_A = \text{det}(A - XI_3) = X^2(X-6)\), on a alors deux sous-espaces propres, \(E_0\) et \(E_6\) et on sait que \(E_0\) est de dimension \(1\), il suffit alors de vérifier que \(E_6\) est bien de dimension \(2\) pour conclure que \(\sum \dim(E_\lambda) = \dim(E)\) et donc que la matrice est diagonalisable.\<

Pour trouver une base de vecteurs propres, il suffit alors de trouver une base de \(E_0, E_6\) et de la concaténer en une base de \(E\).

\chapter*{\chapterstyle{III --- Polynomes d'endomorphismes}}
\addcontentsline{toc}{section}{Polynomes d'endomorphismes}

\chapter*{\chapterstyle{III --- Trigonalisation}}
\addcontentsline{toc}{section}{Trigonalisation}

\chapter*{\chapterstyle{III --- Applications de la réduction}}
\addcontentsline{toc}{section}{Applications de la réduction}

\chapter*{\chapterstyle{III --- Espaces Préhilbertiens Réels}}
\addcontentsline{toc}{section}{Espaces Préhilbertiens Réels}
On appelle \textbf{espace préhilbertien réel} un \(\R\)-espace vectoriel muni d'une forme \textbf{bilinéaire}\footnote[1]{Une application \(\dotproduct{\cdot}{\cdot} : E^2 \longrightarrow \R\) qui est linéaire en chaque variable} qui vérifie:

\customBox{width=7cm}{
   \begin{align*}
      &\textbf{Symétrie} &&\dotproduct{u}{v} = \dotproduct{v}{u} \\
      &\textbf{Définie} &&\dotproduct{u}{u} = 0 \implies u = 0 \\
      &\textbf{Positivité} &&\dotproduct{u}{v} \geq 0
   \end{align*}
}   
On dira alors qu'une telle forme est un \textbf{produit scalaire} sur \(E\).
\subsection*{\subsecstyle{Exemples {:}}}
\begin{itemize}
   \item On définit sur \(\R^n\) le produit scalaire défini par:
   \[
      \dotproduct{u}{v} = \sum_{k=1}^{n} u_kv_k    
   \]
   \item On définit sur \(\mathcal{C}^1(\icc{0}{1}, \R)\) le produit scalaire défini par:
   \[
      \dotproduct{f}{g} = \int_{0}^{1} f(t)g(t) \d t    
   \]
   \item On définit sur \(\R_n[X]\) le produit scalaire défini pour \((x_n)\) \(n + 1\) points fixés par:
   \[
      \dotproduct{P}{Q} = \sum_{k=0}^{n}P(x_k)Q(x_k)
   \]   
   \item On définit sur \(\mathcal{M}_n(\R)\) le produit scalaire défini par:
   \[
      \dotproduct{A}{B} = \text{tr}(A^\top B)
   \]
\end{itemize}

\subsection*{\subsecstyle{Norme induite {:}}}
A partir de la définition d'un tel espace, on alors montrer que \(\dotproduct{\cdot}{\cdot}\) \textbf{induit une norme} sur \(E\) donnée par:
\customBox{width=4.7cm}{
   \(
      \forall u \in E \; ; \; \vectNorm{u}^2 = \dotproduct{u}{u}  
   \)
}
Cela fait donc de \(E\) un espace vectoriel normé.

\subsection*{\subsecstyle{Angle et Orthogonalité{:}}}
A partir d'une telle définition, on peut alors définir \textbf{l'angle non-orienté} \(\theta \in \icc{0}{\pi}\) entre deux vecteurs \(u, v\) par:
\customBox{width=4cm}{
   \[
      \theta := \arccos{\frac{\dotproduct{u}{v}}{\vectNorm{u}\vectNorm{v}}}   
   \]
}
Ce qui nous permet de définir \textbf{l'orthogonalité} de deux vecteurs par la nullité du produit scalaire, en effet \(\theta = \frac{\pi}{2}\) dans cette définition ssi \(\dotproduct{u}{v} = 0\)
\begin{center}
   \textit{L'interprétation de cet "angle" ou de "l'orthogonalité" entre deux vecteurs diffère selon le contexte, elle peut alors signifier une corrélation en probabilité, ou un réel angle géométrique dans \(\R^n\) par exemple.}
\end{center}

\subsection*{\subsecstyle{Inégalité de Cauchy-Schwarz{:}}}
Dans tout espace préhilbertien réel, pour tout \(u, v \in E\) on a l'inégalité\footnote[1]{Trés puissante et permet d'obtenir des majorations dans des cas trés variés.} suivante:
\customBox{width=4cm}{
   \(
      |\dotproduct{u}{v}| \leq \vectNorm{u}\,\vectNorm{v}
   \)
}
Avec cas d'égalité quand \(u, v\) sont \textbf{liés}.

\subsection*{\subsecstyle{Formules géométriques{:}}}
Dans tout espace préhilbertien réel, pour tout \(u, v \in E\) on a les identités suivantes:
\begin{itemize}
   \item \textbf{Identité de polarisation :} \(\vectNorm{x + y}^2 = \vectNorm{x}^2 + \vectNorm{y}^2 + 2 \dotproduct{x}{y}\).
   \item \textbf{Identité du parallélogramme :} \(\vectNorm{x + y}^2 + \vectNorm{x - y}^2 = 2\vectNorm{x}^2 + 2\vectNorm{y}^2\).
\end{itemize}
Enfin, on a le \textbf{théorème de Pythagore}:
\[
   x \perp y \Longleftrightarrow \vectNorm{x + y}^2 = \vectNorm{x}^2 + \vectNorm{y}^2
\]

\subsection*{\subsecstyle{Orthogonal d'une partie{:}}}
Soit \(F\) une partie de \(E\), alors on peut définir l'orthogonal de \(F\) comme \textbf{le sous-espace vectoriel} défini par:
\customBox{width=6cm}{
   \(F^\perp := \bigl\{ u \in E \; ; \; \forall v \in F \, , \, u \perp v  \bigl\}\) 
}
\begin{center}
   \textit{L'orthogonal d'une partie est l'ensemble des vecteurs qui sont orthogonaux à tout élément de cette partie.}
\end{center}
Dans le cas courant ou on connait une base \((e_1, \ldots, e_n)\) de \(F\), on a la caractérisation suivante:
\customBox{width=7cm}{
   \(F^\perp := \bigl\{ u \in E \; ; \; \forall i \in \inticc{1}{n} \, , \, u \perp e_i  \bigl\}\)
}
On peut alors montrer les propiétés suivantes sur l'orthogonal:
\begin{itemize}
   \item \(F \subseteq {(F^\perp)}^\perp\)
   \item \(E^\perp = \{ 0_E \}\)
   \item \(A \subseteq B \implies B^\perp \subseteq A^\perp\)
   \item \((A \cup B)^\perp =  A^\perp \cap B^\perp\)
\end{itemize}
Enfin on peut aussi montrer la propriété fondamentale suivante:
\begin{center}
   \textbf{Une partie et son orthogonal sont toujours en somme directe.}
\end{center}
Attention dans un espace préhilbertien quelconque, ils ne sont pas toujours supplémentaires !

\chapter*{\chapterstyle{III --- Espaces Euclidiens}}
\addcontentsline{toc}{section}{Espaces Euclidiens}
On appelle \textbf{espace euclidien} tout espace préhilbertien réel \textbf{de dimension finie}. Dans toute la suite on prendra \((e_1, \ldots, e_n)\) une base de \(E\).

\subsection*{\subsecstyle{Forme analytique et matricielle du produit scalaire{:}}}
En représentant les vecteurs \(x, y \in E\) dans la base, on peut alors développer leur produit scalaire par bilinéarité pour obtenir:
\customBox{width=5.5cm}{
   \[
      \dotproduct{x}{y} = \sum_{i = 1}^{n}\sum_{j = 1}^{n}x_iy_j \dotproduct{e_i}{e_j}   
   \]
}
En particulier on peut alors remarquer que le produit scalaire est parfaitement déterminé par la donnée des \(n^2\) produits scalaires des vecteurs de base. Il est donc représenter matriciellement le produit scalaire par la matrice
\(
   A = (a_{i,j}) \text{ telle que } a_{i, j} = \dotproduct{e_i}{e_j}   
\)\<

Et donc on obtient en identifiant \(\R\) à \(\mathcal{M}_{1, 1}(\R)\):
\customBox{width = 3.7cm}{
   \(
      \dotproduct{x}{y} = X^\top A Y
   \)
}
En particulier, dans une base \textbf{orthonormée}, on a \(\dotproduct{x}{e_i} = x_i\) et donc:
\[
   x = \sum_{k=1}^{n} \dotproduct{x}{e_i}e_i 
\]
\subsection*{\subsecstyle{Propriétés en dimension finie{:}}}
En dimension finie, pour tout sous-espace vectoriel \(F\) de \(E\) on a le théorème fondamental suivant:
\customBox{width = 5cm}{
   \(
      E = F \oplus F^\perp    
   \)
}
\begin{center}
   \textit{Tout sous-espace admet un unique supplémentaire orthogonal en dimension finie.}
\end{center}
On peut alors en déduire qu'en dimension finie on a:
\[
   (F^\perp)^\perp = F   
\]

\subsection*{\subsecstyle{Projection orthogonale{:}}}
L'existence d'une unique décomposition nous permet alors de définir \textbf{la projection} sur \(F\) de direction \(F^\perp\) par:
\[
   \text{proj}_F : x = x_F + x_{F^\perp} \mapsto x_F   
\]
En particulier, on en déduit par l'unicité de la décomposition que \(\text{proj}_F(x)\) est \textbf{l'unique vecteur} de \(F\) qui vérifie:
\customBox{width=4cm}{
   \(x - \text{proj}_F(x) \in F^\perp\)
}
Le projeté orthogonal a une signification géométrique importante, en effet on a:
\customBox{width=6cm}{
   \(\vectNorm{x - \text{proj}_F(x)} = \min_{y\in F}(\vectNorm{x - y})\)
}
\begin{center}
   \textit{C'est le vecteur "le plus proche" de \(F\) au sens du produit scalaire utilisé.}
\end{center}
\pagebreak

\subsection*{\subsecstyle{Calcul de projeté orthogonal{:}}}
On considère un sous-espace \(F\) de bases \((e_1, \ldots, e_p)\) et \(x \in E\), on sait d'aprés les propriétés précédentes que \(\text{proj}_F(x)\) est l'unique vecteur de \(F\) tel que \(x - \text{proj}_F(x) \in F^\perp\), ce qui est équivalent à dire que:
\begin{align*}
   &\forall j \in \inticc{1}{p} \; ; \; \dotproduct{x - \text{proj}_F(x)}{e_j} = 0 \Longleftrightarrow \\
   &\forall j \in \inticc{1}{p} \; ; \; \dotproduct{\text{proj}_F(x)}{e_j} = \dotproduct{x}{e_j}
\end{align*}
On raisonne alors par coefficient indeterminés avec l'écriture de \(\text{proj}_F(x) = \sum_{i=1}^{n} \alpha_i e_i\) dans la base de \(F\) pour obtenir l'expression suivante:
\[
   \forall j \in \inticc{1}{p} \; ; \; \sum_{i=1}^{n}\alpha_k \dotproduct{e_i}{e_j} = \dotproduct{x}{e_i}
\] 
Enfin on obtient alors \textbf{le système des équations normales}:
\[
   \begin{cases}
      \alpha_1 \dotproduct{e_1}{e_1} + \alpha_2 \dotproduct{e_2}{e_1} + \ldots + \alpha_p \dotproduct{e_p}{e_1} = \dotproduct{x}{e_1}\\
      \vdots \\
      \alpha_1 \dotproduct{e_1}{e_p} + \alpha_2 \dotproduct{e_2}{e_p} + \ldots + \alpha_p \dotproduct{e_p}{e_p} = \dotproduct{x}{e_p}
   \end{cases} 
\]
Ou de manière équivalente pour \(M\) la matrice du produit scalaire dans la base de \((e_1, \ldots, e_p)\):
\[
   MY = 
   \begin{bmatrix}
      \dotproduct{x}{e_1}\\
      \vdots \\
      \dotproduct{x}{e_p}
   \end{bmatrix}
\]
Dans le cas d'une base \textbf{orthogonale}, le système ci-dessus est beaucoup plus simple, en effet presque tout les produits scalaires sont nuls, et on obtient un système \textbf{diagonal} et donc dans ce cas précis, le projeté orthogonal s'obtient simplement par la formule:
\customBox{width=5cm}{
   \[
      \text{proj}_F(x) = \sum_{k = 1}^{p} \frac{\dotproduct{x}{e_i}}{\dotproduct{e_i}{e_i}}e_i     
   \]
}
Finalement, une remarque importante permet de comprendre la projection sur un sous espace doté d'une base orthogonale, en effet:
\begin{center}
   \textit{Projeter un vecteur sur un sous-espace revient à ajouter les projetés de ce vecteur sur les vecteurs de la base du sous-espace.}
\end{center}
\underline{Exemple:} Le projeté de \(x = (1, 2, 3)\) sur le plan \(\text{Vect}((1, 0, 0), (0, 1, 0))\) est donné par \(\text{proj}_{(1, 0, 0)}(x) + \text{proj}_{(0, 1, 0)}(x)\)
\subsection*{\subsecstyle{Procédé de Gramm-Schmidt{:}}}
Soit \((e_1, \ldots, e_n)\) une base de \(E\), on cherche alors à élaborer un procédé permettant \textbf{d'orthogonaliser cette base} en une base \((\epsilon_1, \ldots, \epsilon_n)\), on pose \(\epsilon_1 = e_1\) et \(H_i = \text{Vect}(e_1, \ldots, e_i)\) et on définit par récurrence:
\customBox{width=12cm}{
   \[
      \epsilon_i = e_i - \text{proj}_{H_{i-1}}(e_i) = e_i - \left(\sum_{k=0}^{i - 1}\text{proj}_{\epsilon_k}(e_k)\right) = e_i - \left(\sum_{k=0}^{i - 1}\frac{\dotproduct{e_i}{\epsilon_i}}{\dotproduct{\epsilon_i}{\epsilon_i}}\epsilon_i \right)
   \]
}

\begin{center}
   \textit{Moralement, on "redresse" chaque vecteur de la base initiale en lui retirant son défaut d'orthogonalité représenté par sa projection sur le sous-espace précédent.}
\end{center}